\subsection{Robustness to the adversary model}

The numbers above are measured against an adversary that reads the responses it
receives, which is what real tooling does: an injection attempt is only informative
if you look at what comes back. An earlier version of our round-2 attacker did not,
and a probe that nobody reads cannot work by construction, so that model made the
probe unreachable rather than ineffective. We report both conditions rather than
replacing one with the other, because the blind attacker is the honest lower bound
for a reader who doubts our curiosity model.

\begin{table}[t]
\centering
\small
\caption{The same frozen model against two adversaries. Only the attacker differs;
the model, the seeds and the traffic volumes are identical. The blind condition is
the floor, not a competing result.}
\label{tab:adversary}
\begin{tabular}{lcc}
\toprule
 & blind adversary & reading adversary \\
\midrule
B2 passive recall & 0.915 & 0.917 \\
B4 full recall & 0.933 & \textbf{0.951} \\
B4 $-$ B2 & +0.018 & \textbf{+0.034} \\
Causal effect (holdout) & +0.029 & \textbf{+0.051} \\
Benign sessions diverted & 0 & 0 \\
\bottomrule
\end{tabular}
\end{table}

Two things are worth reading off Table~\ref{tab:adversary}. The probe's advantage
roughly doubles once the adversary behaves realistically, so the blind figure
understates it. And B2, which has no probe at all, barely moves (0.915 to 0.917):
making the attacker read responses does not make attacks easier to catch in
general, it only makes the probe reachable. That control is what lets us attribute
the difference to the probe rather than to a friendlier adversary.

